Graph traversal over large sparse datasets is a latency-critical primitive underlying
financial fraud detection, graph databases, genomics assembly, and real-time recommendation.
General-purpose processors fail this workload: CPU caches thrash on random pointer chains,
and GPU kernel dispatch overhead precludes deterministic sub-microsecond response.
A companion analysis~\cite{fritz2026rmwwall} establishes that graph BFS achieves only
$\eta \approx 40\%$ bandwidth efficiency on HBM4 due to row-activation scheduling
constraints, and that this gap is not closed by increasing data rate.

We present \textit{BFS-HBM}: a fully open-source, synthesizable SystemVerilog
microarchitecture for BFS traversal of CSR-encoded graphs stored in High-Bandwidth Memory.
The design comprises a 12-state traversal FSM, an AXI4-compliant read master, a
FIFO-based frontier queue, and a 2-cycle BRAM read-modify-write visited bitmap.
Graph data is laid out in Compressed Sparse Row (CSR) format; a single 256-bit
HBM beat delivers eight 32-bit neighbor IDs, and both endpoints of the CSR row pointer
are fetched in one transaction, eliminating a second memory round trip per vertex.

Cycle-accurate RTL simulation at a modeled HBM latency of 20 cycles verifies
correct BFS discovery ordering and level assignment across all 8 vertices of an
undirected test graph in 514 cycles at 250~MHz. Throughput analysis shows that
the dominant bottleneck is the 2-cycle serial visited-bitmap RMW loop, which
limits steady-state throughput to one edge per 2 cycles — consistent with the
$\eta \approx 40\%$ bandwidth efficiency of the underlying HBM interface.

The RTL is parameterized by vertex count ($2^{\mathit{VERTEX\_W}}$ vertices, default
$2^{20} = 1\text{M}$), AXI data width, and FIFO depth, and targets Xilinx UltraScale+
HBM devices (Alveo U280) for FPGA validation and TSMC N7 / GlobalFoundries 12LP for
ASIC tapeout. All source, testbench, and synthesis scripts are released under the
MIT license.
